\scanpage{16}
\mspage{\handwritten{30}}

\deleted{\deleted{pourra alors}\inserted{\handwritten{peut donc}} \inserted{dans 3.9.} resper l'hypothèse \deleted{dans 3.9.} avec : ``ou \deleted{\unsure{$\mathrm B$|$\mathrm B_0$} \unsure{est un anneau un corp|est un anneau un corps}} $\mathrm r(\mathrm B)=\mathrm r(\mathrm A)\mathrm B$'' \deleted{(car grâce au résultat de et la}}

\uline{Corollaire} \deleted{\unsure{3|5}}\handwritten{4}.10. Soi\deleted{t}ent $\mathrm A$ un anneau local noethérien, de corps résiduel $k$, $\mathrm B_{\unsure{0|o}}$ une $k$-algèbre qui soit un anneau local noethérien, et soit une $k$-algèbre de Cohen i.e. une $k$-algèbre formel\deleted{\unsure{h|l}}lement simple \deleted{pour sa} et complète pour sa topologie ordinaire. Supposons de plus l'extension résiduelle $\mathrm K/k$ séparable (ce qui n'est donc pas une restriction si $\mathrm B_0=\mathrm K$\deleted{\unsure{=k|,}}, mais en est évidemment une en général). Alors il existe une $\mathrm A$-algèbre topologique de Cohen \inserted{$\mathrm B$} qui se réduit suivant $\mathrm B_0$, $\mathrm B$ est unique à isomorphisme près en vertu de \deleted{3}\handwritten{4}.8. , et c'est un anneau noethérien local complet comme il se doit.

En vertu du théorème de structure \deleted{\unsure{2}}\handwritten{3}.4. on a $\mathrm B_0\handwritten{\simeq}\mathrm K\handwritten{[\![}\mathrm T_1,\ldotp\ldotp,\mathrm T_n\handwritten{]\!]}$, et on est ramené, grâce au critère \deleted{3}\handwritten{4}.9. 2) et la transitivité de la platitude, aux deux c\deleted{\unsure{q}}as types: 1\textsuperscript{\uline{o}}) $\mathrm B_0=\mathrm K$ 2\textsuperscript{\uline{o}}) $\mathrm B_0=k\handwritten{[\![}\mathrm T_1,\ldotp\ldotp,\mathrm T_n\handwritten{]\!]}$. Dans le premier cas, on utilise le résultat général de $0_{\mathrm{III}}$\inserted{\handwritten{10.3.1.}}, et dans le deuxième cas, il suffit de prendre $\mathrm B=\handwritten{\widehat{\mathrm A}[\![}\mathrm T_1,\ldots,\mathrm T_n\handwritten{]\!]}$.

\deleted{Théorème 3.1} \uline{Corollaire} \deleted{3}\handwritten{4}.11. Sous les conditions de 3.9. , \inserted{\handwritten{si $\unsure{K|k}(B)/k(A)$ de mult. rad. finie,}} les conditions 1) et 2) équivalent aussi à: 3) $\mathrm B$ est complet et formellement simple sur $\mathrm A$

En effet, on sait déjà que \deleted{c'est nécessaire,}\inserted{1)$\handwritten{\Rightarrow}$3)}, reste à prouver \deleted{la} 3)$\handwritten{\Rightarrow}$1) \deleted{suffisance} ou encore 3)$\handwritten{\Rightarrow}$2), i.e. on est ramené à prouver simplement que $\mathrm B$ est plat sur $\mathrm A$. Or\deleted{\unsure{n}} faisant une extension de la base $\mathrm A\handwritten{\to}\mathrm A'$, $\mathrm A'$ local libre \inserted{de type fini} sur $\mathrm A$ à extension radicielle convenable (cf $0_{\mathrm{III}}$), on trouve un $\mathrm B'$ \inserted{local} sur $\mathrm A'$ qui sera encore formellement simple sur $\mathrm A'$, et de plus à extension résiduelle \deleted{triviale} séparable\handwritten{$/^{\circ}$}. On dispose alors du théorème d'existence \deleted{3}\handwritten{4}.10. qui nous garantit, grâce à \deleted{3}\handwritten{4}.8. (équiv de 0) et 1) ) que $\mathrm B'$ est Cohen sur $\mathrm A'$donc plat sur $\mathrm A'$. Donc $\mathrm B$ est plat sur $\mathrm A$.

\margin{\handwritten{($/^{\circ}$ réduction inutile dans le cas ``\unsure{vespé|visé|vrai}'' du NB)}}
% archive sheet number: 15
